
\documentclass{article}
\usepackage{colortbl}
\usepackage{makecell}
\usepackage{multirow}
\usepackage{supertabular}

\begin{document}

\newcounter{utterance}

\centering \large Interaction Transcript for game `cladder', experiment `full\_v1.5\_default', episode 9980 with reflexion{-}llama{-}3.1{-}8b{-}instruct.
\vspace{24pt}

{ \footnotesize  \setcounter{utterance}{1}
\setlength{\tabcolsep}{0pt}
\begin{supertabular}{c@{$\;$}|p{.15\linewidth}@{}p{.15\linewidth}p{.15\linewidth}p{.15\linewidth}p{.15\linewidth}p{.15\linewidth}}
    \# & \multicolumn{2}{c}{Player} && \multicolumn{2}{c}{Game Master} \\
    \hline

    \theutterance \stepcounter{utterance}  
    & & & \multicolumn{4}{p{0.6\linewidth}}{
        \cellcolor[rgb]{0.9,0.9,0.9}{
            \makecell[{{p{\linewidth}}}]{
                \texttt{\tiny{[P1$\langle$GM]}}
                \texttt{You are an expert in causal inference. The following question is not a typical commonsense query, but rather a meticulously designed question created by a professor specializing in causal inference, intended to assess the students' mastery of the course content.} \\
\texttt{Imagine a self{-}contained, hypothetical world with only the following conditions, and without any unmentioned factors or causal relationships: Hwax has a direct effect on pexu and rukz. Pexu has a direct effect on rukz. We know that hwax causes pexu. hwax and pexu causes rukz. We observed an individual is not hwax. Would an individual is not rukz if not pexu instead of pexu?} \\
\texttt{Start your answer with 'Yes' or 'No', followed by additional reasoning or evidence to support your explanation.} \\
            }
        }
    }
     \\ \\

    \theutterance \stepcounter{utterance}  
    & \multicolumn{4}{p{0.6\linewidth}}{
        \cellcolor[rgb]{0.561,0.737,0.561}{
            \makecell[{{p{\linewidth}}}]{
                \texttt{\tiny{[P1$\rangle$GM]}}
                \texttt{No.} \\
\\ 
\texttt{To answer this question, we need to apply the principles of causal inference, specifically the concept of causal chains and the backdoor criterion.} \\
\\ 
\texttt{Given the conditions:} \\
\\ 
\texttt{1. Hwax causes pexu.} \\
\texttt{2. Hwax causes rukz.} \\
\texttt{3. Pexu causes rukz.} \\
\\ 
\texttt{We can represent these relationships using a causal diagram:} \\
\\ 
\texttt{Hwax {-}> Pexu} \\
\texttt{Hwax {-}> Rukz} \\
\texttt{Pexu {-}> Rukz} \\
\\ 
\texttt{Now, we are given that the individual is not exposed to Hwax (i.e., Hwax = 0). We want to know whether the individual would be rukz if they were not exposed to Pexu instead (i.e., Pexu = 0).} \\
\\ 
\texttt{To answer this question, we need to consider the backdoor paths between Hwax and Rukz. In this case, there is a backdoor path through Pexu: Hwax {-}> Pexu {-}> Rukz.} \\
\\ 
\texttt{However, since we are conditioning on Pexu = 0, we need to consider the effect of Hwax on Rukz through Pexu. If Pexu = 0, the effect of Hwax on Rukz through Pexu is blocked, and the only remaining effect of Hwax on Rukz is direct (Hwax {-}> Rukz).} \\
\\ 
\texttt{Since we know that Hwax causes rukz, and the individual is not exposed to Hwax, we can conclude that the individual would still be rukz even if they were not exposed to Pexu.} \\
\\ 
\texttt{Therefore, the answer is No: the individual would still be rukz if not pexu instead of pexu.} \\
            }
        }
    }
    & & \\ \\

    \theutterance \stepcounter{utterance}  
    & & & \multicolumn{2}{p{0.3\linewidth}}{
        \cellcolor[rgb]{0.9,0.9,0.9}{
            \makecell[{{p{\linewidth}}}]{
                \texttt{\tiny{[GM$|$GM]}}
                \texttt{No} \\
            }
        }
    }
    & & \\ \\

    \theutterance \stepcounter{utterance}  
    & & & \multicolumn{2}{p{0.3\linewidth}}{
        \cellcolor[rgb]{0.9,0.9,0.9}{
            \makecell[{{p{\linewidth}}}]{
                \texttt{\tiny{[GM$|$GM]}}
                \texttt{yes} \\
            }
        }
    }
    & & \\ \\

    \theutterance \stepcounter{utterance}  
    & & & \multicolumn{2}{p{0.3\linewidth}}{
        \cellcolor[rgb]{0.9,0.9,0.9}{
            \makecell[{{p{\linewidth}}}]{
                \texttt{\tiny{[GM$|$GM]}}
                \texttt{game\_result = LOSE} \\
            }
        }
    }
    & & \\ \\

    \theutterance \stepcounter{utterance}  
    & & \multicolumn{4}{p{0.6\linewidth}}{
        \cellcolor[rgb]{0.95,0.95,0.95}{
            \makecell[{{p{\linewidth}}}]{
                \texttt{\tiny{[P1$\leftrightarrow$P]}}
                \texttt{Post{-}episode Reflexion (not sent to GM):} \\
\texttt{New plan: } \\
\\ 
\texttt{In my previous attempt, I correctly identified the causal relationships between Hwax, Pexu, and Rukz, and applied the backdoor criterion to determine the effect of Hwax on Rukz. However, I failed to consider the implications of conditioning on Pexu = 0 on the causal chain Hwax {-}> Pexu {-}> Rukz. } \\
\\ 
\texttt{To improve my strategy, I will pay closer attention to the conditioning statements and their effects on the causal relationships. In this case, conditioning on Pexu = 0 blocks the backdoor path through Pexu, but I should have also considered the direct effect of Hwax on Rukz (Hwax {-}> Rukz). } \\
\\ 
\texttt{Next time, I will carefully analyze the causal diagram and the conditioning statements to ensure that I account for all possible causal paths and their implications. I will also double{-}check my reasoning to avoid overlooking critical details, such as the direct effect of Hwax on Rukz in this case.} \\
            }
        }
    }
    & \\ \\

\end{supertabular}
}

\end{document}
